\documentclass[../main.tex]{subfiles}
\begin{document}
%==============Tờ số 1==================
\begin{center}
    \begin{overpic}[height=1.3\pagewidth]{../../../../Image/Book?/chaptitle/intro1.png}
    \put(50, 50){               % x = 50 (mặc định trong quyển này), y = 50
        \begin{tcolorbox}[
            colframe=oldpaper,  % Màu viền
            boxrule=1pt,      % Độ dày viền
            arc=0mm,         % Góc vuông (không bo tròn)
            enhanced jigsaw,
            opacityback = 0,
            opacityframe = 0,
            left=5pt,        % Khoảng cách giữa nội dung và viền trái
            right=5pt,       % Khoảng cách giữa nội dung và viền phải
            top=5pt,         % Khoảng cách giữa nội dung và viền trên
            bottom=5pt,      % Khoảng cách giữa nội dung và viền dưới
            width = 13cm     % (mặc định)
        ]
        {\color{oldpaper}
        Tháng Bảy, năm 2018.
        
        \vspace{10pt}
    
        Tôi là Haragen Kyuen, và có cô em gái song sinh là Haragen Saikyou. Tôi còn chẳng biết liệu đó có thật sự là tên chúng tôi không, vì\dots\ tôi còn chẳng nhớ rõ mình là ai nữa.
    
        \vspace{5pt}
    
        Chúng tôi bị dịch chuyển về quá khứ, và được trải nghiệm lại những gì đã xảy ra vào ngày hôm đó. Chúng tôi không rõ tại sao lại bị rơi vào mớ bòng bong như vậy.
    
        \vspace{5pt}
    
        Ngày đầu tiên được học tại đây, một ngôi trường trung học ở ngoại thành có tên là Trường Trung học Aogami.
    
        \vspace{5pt}
    
        Mọi chuyện bắt đầu từ việc sau khi đã kết thúc tiết học đầu tiên, và bắt đầu giờ nghỉ giải lao giữa giờ\dots
    
        \vspace{5pt}
    
        Thoạt nhìn, đây có lẽ chỉ là một ngôi trường rất bình thường, thế nhưng\dots
    
        \vspace{5pt}
    
        Tôi có nghe qua một vài câu chuyện ma tại ngôi trường Aogami này. Một trong số đó là:
    
        \vspace{5pt}
    
        ``\textbf{Nếu bạn không đối xử tốt với một học sinh mới chuyển trường, những điều bất hạnh sẽ xảy ra.}''
    
        \vspace{5pt}
    
        Và rồi\dots}
        \end{tcolorbox}
    }
    \end{overpic}
\end{center}
%==============Tờ số 2==================
\newpage
\thispagestyle{empty}
\begin{center}
    \begin{overpic}[height=1.3\pagewidth]{../../../../Image/Book?/chaptitle/intro0.png}
    \put(50, 700){
        \begin{tcolorbox}[
            colframe=oldpaper,  % Màu viền
            boxrule=1pt,      % Độ dày viền
            arc=0mm,         % Góc vuông (không bo tròn)
            enhanced jigsaw,
            opacityback = 0,
            opacityframe = 0,
            left=5pt,        % Khoảng cách giữa nội dung và viền trái
            right=5pt,       % Khoảng cách giữa nội dung và viền phải
            top=5pt,         % Khoảng cách giữa nội dung và viền trên
            bottom=5pt,      % Khoảng cách giữa nội dung và viền dưới
            width = 13cm     % (mặc định)
        ]
        {\color{oldpaper}
        Ngày X Tháng Bảy, năm 2018.
    
        \vspace{10pt}
    
        Một vụ tai nạn thảm khốc đã xảy ra tại Trường Trung học Aogami khiến cho cả ngôi trường đổ sập.
    
        \vspace{5pt}
    
        Cảnh sát địa phương và Sở Cứu hỏa của thị trấn đã bắt tay điều tra vụ việc, tuy nhiên, họ đã bắt gặp nhiều hiện tượng kỳ bí tại hiện trường, và nguyên nhân vụ việc vẫn chưa được làm rõ.
    
        \vspace{5pt}
    
        Vụ việc này vẫn còn rất nhiều câu hỏi bỏ ngỏ, trong đó: 
        }
        \end{tcolorbox}}
    \put(50, 440){
        \begin{tcolorbox}[
            colframe=oldpaper,  % Màu viền
            boxrule=3pt,      % Độ dày viền
            arc=0mm,         % Góc vuông (không bo tròn)
            enhanced jigsaw,
            opacityback = 0,
            left=5pt,        % Khoảng cách giữa nội dung và viền trái
            right=5pt,       % Khoảng cách giữa nội dung và viền phải
            top=5pt,         % Khoảng cách giữa nội dung và viền trên
            bottom=5pt,      % Khoảng cách giữa nội dung và viền dưới
            width = 13cm     % (mặc định)
        ]
        {\color{oldpaper}
        \begin{center}
            {\arial \textbf{Những điểm bất thường}}
        \end{center}
    
        \vspace{10pt}
    
        \begin{itemize}
            \item Người ta cho rằng có một số học sinh và nhân viên liên quan tới vụ tai nạn, nhưng vẫn chưa biết được \emph{liệu họ có an toàn hay không}.
            \item Mặc dù những người dân địa phương cho rằng \emph{một trận động đất lớn đã xảy ra cùng thời điểm với vụ tai nạn trên}, nhưng Cơ quan Khí tượng Quốc gia lại không có báo cáo ghi nhận điều đó.
            \item Dạnh sách học sinh được phía nhà trường đưa ra \emph{không khớp với số lượng người bị mất tích được báo cáo.}
        \end{itemize}}
        \end{tcolorbox}}
    \put(50, 380){
        \begin{tcolorbox}[
            colframe=oldpaper,  % Màu viền
            boxrule=1pt,      % Độ dày viền
            arc=0mm,         % Góc vuông (không bo tròn)
            enhanced jigsaw,
            opacityback = 0,
            opacityframe = 0,
            left=5pt,        % Khoảng cách giữa nội dung và viền trái
            right=5pt,       % Khoảng cách giữa nội dung và viền phải
            top=5pt,         % Khoảng cách giữa nội dung và viền trên
            bottom=5pt,      % Khoảng cách giữa nội dung và viền dưới
            width = 13cm     % (mặc định)
        ]
        {\color{oldpaper}
        Cảnh sát đã kết luận nguyên nhân vụ tai nạn này xảy ra là do một vụ nổ từ rò rỉ khí ga.}
        \end{tcolorbox}
    }
    \end{overpic}
\end{center}
\end{document}